<<========================== UPDATE ==========================>
Tius : (28/10/2026)
        1. Perbaiki tampilan pada dashboard
        2. Hapus file tidak penting (scss)
        3. data sementara dari pendaftar.json
        4. animasi dashboard menggunakan js dan dipisahkan
        5. Memperbarui Head.php

Tius : (01/11/2026)
        1. Update dashboard
        2. Update Data Mahasiswa
        3. Update Database
        4. Update Form Tambah Mahasiswa
        5. Hapus Data Json
        6. Update struktur syntax

Ananta : (02/11/2026)
        log :
        1. Update struktur Sidebar (Tema ungu, collapse feature, nested dropdown)
        2. Update Modul Data Mahasiswa (Pemisahan Camaba & Mahasiswa Resmi, Live Search, Paginasi)
        3. Update Database (Tabel prodi, update ENUM status, training data uji)
        4. Update Form Tambah Mahasiswa (Validasi duplikat NIM/NIK, kalender TTL, dropdown prodi)
        5. Integrasi Modul Klasifikasi Naive Bayes (Otomatisasi status Menunggu -> Diterima/Ditolak)
        6. Perbaikan bug Notice Session & perbaikan layout tabel

        detail :
        1. Modul Antarmuka & Sidebar (UI/UX)
                a. Pembaruan Tema: Menyelaraskan seluruh warna sistem menjadi identitas Ungu (Purple Palette) yang dikombinasikan dengan latar belakang gray-900 yang elegan.
                b. Integrasi Logo: Memasukkan image.png pada bagian atas sidebar bersanding dengan teks "SPK.PMB".
                c. Fitur Collapse Sidebar: Menambahkan tombol interaktif untuk melipat (collapse) dan membuka sidebar agar ruang kerja tabel menjadi lebih luas dan responsif.
                d. Restrukturisasi Menu (Nested Dropdown): Memasukkan menu "Data Mahasiswa" ke dalam "Master Data" dengan dropdown bertingkat (nested) berisi "Data Camaba" dan "Mahasiswa Resmi". 
                Memperbaiki bug event bubbling saat menu diklik.
        2. Modul Data Mahasiswa & Camaba
                a. Pemisahan Tampilan: Membedakan rute halaman "Data Camaba" (menampilkan semua pendaftar dengan filter status) dan "Mahasiswa Resmi" (eksklusif hanya menampilkan status 'Diterima' tanpa tombol tambah data).
                b. Pemisahan Tampilan: Membedakan rute halaman "Data Camaba" (menampilkan semua pendaftar dengan filter status) dan "Mahasiswa Resmi" (eksklusif hanya menampilkan status 'Diterima' tanpa tombol tambah data).
                c. Paginasi & Limit Dinamis: Mengembalikan desain kotak form asli dan memindahkan opsi dropdown "Batasi Data" ke bagian header tabel. Menambahkan sistem Live Search dan paginasi data.
        3. Modul Form Pendaftaran (Validasi Ketat)
                a. Validasi Karakter: NIM dibatasi maksimal 12 digit dan NIK 16 digit.
                b. Pencegahan Duplikat: Menambahkan proteksi validasi di MahasiswaModel.php agar tidak ada NIM atau NIK yang sama di dalam database (baik saat Tambah maupun Edit).
                c. Otomatisasi Status: Menghilangkan pilihan "Status Kelulusan" manual. Pendaftar baru secara mutlak akan berstatus "Menunggu" sebelum dianalisis sistem.
                d. Integrasi Kalender & Prodi: Memisahkan input Tempat Lahir (teks) dan Tanggal Lahir (Date Picker Browser). Mengubah isian Program Studi Tujuan dari input teks menjadi dropdown yang memanggil data asli dari tabel.
        4. Modul Analisis Naive Bayes (Inti SPK)
                a. Pembuatan Arsitektur MVC: Membangun AnalisisModel.php, AnalisisController.php, dan analisis.php dari nol sebagai pengganti konsep WP yang keliru.
                b. Algoritma Probabilitas: Menulis rumusan perhitungan Prior Probability, Conditional Probability, dan Laplace Smoothing murni untuk mengklasifikasi status pendaftar yang "Menunggu".
                c. Perbaikan Pembacaan Ciri: Algoritma dioptimalkan untuk membaca id_sub (Opsi Pilihan Kriteria) alih-alih angka desimal mentah.
        5. Pembaruan Database (SQL)
                a. Pembuatan tabel prodi berisi 8 daftar jurusan Universitas Triatma Mulya.
                b. Pembaruan struktur ENUM pada kolom status tabel mahasiswa menjadi ('Menunggu', 'Diterima', 'Ditolak').
                c. Pembaruan struktur ENUM pada kolom status tabel mahasiswa menjadi ('Menunggu', 'Diterima', 'Ditolak').